\section{问题分析}
\subsection{问题1的分析}
问题一需要在观测要素不同的条件下建立具有共同输出含义的两个模型。模型 a 同时利用热力廓线和风场，能够计算层结稳定度及动力切变；模型 b 的输入中缺少温度信息，需要从谱宽、风速梯度和历史变化中提取与模型 a 有关的统计特征。两种模型的区别因此体现在信息来源，而非输出单位。

建模时先固定模型 a 的计算规则，再在训练时间块上标定模型 b。归一化参数、缺测处理阈值和回归系数均由训练部分估计，测试部分用于比较整体误差与分高度误差。这样可以分别讨论参数拟合效果和跨时段变化，避免将相邻时刻的重复信息当作独立验证。
\subsection{问题2的分析}
问题二的核心是利用不均匀分布的观测恢复规则网格上的场。地面站约束近地面，风廓线雷达约束站点上方的垂直结构，天气雷达提供扫描区域的信息。首先按设备定义构造可比较的指标，再通过观测算子描述采样体积与网格值的关系。空间权重由距离、误差和时效共同决定，平滑项用于约束未直接观测的位置。

在结果检验中，应将整个站点或连续空间块留出后重新求解。若只移除一个点而保留其附近的同设备重复观测，误差会更多反映插值的局地连续性。本文因此同时报告留站误差、观测覆盖和参数扰动下的场变化。
\subsection{问题3的分析}
问题三将监测转化为预报，再将预报转化为路径代价。模型 d 包含数值模式的未来信息，模型 e 依靠历史场的演化规律，两者的输入边界不同，需要分别验证。采用逐次推进的预报起点，比较同一时效上的误差，最终以截至 05:00 的信息生成题目要求的未来场。

航路规划中，节点应包含位置和到达时间，使风险读取与飞行过程对应。图搜索的最优性取决于非负边权、允许区域和启发函数下界。对于只有湍流暴露量的目标，先采用零启发函数建立 Dijkstra 基线，再在已知单位距离风险下界时使用 A*。
\auxfigure{analysis}{问题难点与对应建模环节}{fig:analysis}
\begin{table}[H]\centering\auxtablecaption{关键难点、处理方法与检验方式}
\begin{tabularx}{\linewidth}{YYY}\toprule
关键难点 & 处理方法 & 检验方式\\\midrule
热力信息缺失 & 雷达特征标定模型 b & 时间块与高度分层误差\\
观测尺度不同 & 统一指标和观测算子 & 留站重构、设备消融\\
短历史序列外推 & 非线性回归与平流基线 & 滚动预报起点回报\\
动态风险下寻路 & 时空图与累计边权 & 小图穷举、Dijkstra 对照\\\bottomrule
\end{tabularx}\end{table}
